\chapter{Theoretical Background }

In this chapter, we will look at the RISC-V ISA, analyze the main elements of developing 
a CPU with SpinalHDL \cite{SpinalHDL2018} on the case of \textbf{VexiiRiscV},
look at the litex ecosystem and framework,
how custom extensions are ususally implemented, and how software is run on the SoCs produced with such
methods.

\section{The RISC-V Architecture and VexiiRiscv Core}

\subsection{Overview of the RISC-V ISA}

The RISC-V instruction set architecture (ISA) \cite{RISC-VInternational2025} is an open standard  designed with simplicity, extensibility, and long-term sustainability in mind.
Unlike proprietary architectures such as x86 or ARM, RISC-V is developed under an
open governance model and can be freely implemented without licensing constraints.

A central design principle of RISC-V is modularity.
The base ISA provides only essential functionality,
while optional extensions add specialized capabilities.
Common extension subsets include integer multiplication (\textbf{IM}),
atomic operations (\textbf{A}),
single and double precision floating point arithmetic (\textbf{F/D})
and compressed 16-bit instructions (\textbf{C}).

Typical Linux-capable systems use configurations such as RV64IMAC \cite{Rakhmatov2023Survey},
which includes general-purpose, compressed, floating-point,
and atomic extensions, however, for our
implementation, we use RV32IMA. \cite{LinuxOnLiteXGitHub}

Importantly, RISC-V reserves opcode space for custom extensions.
These custom instruction fields enable designers to introduce application-specific
accelerators directly into the processor pipeline without
conflicting with standard instructions.
This extensibility forms the basis for the custom CXU interface described later in this chapter.

\subsection{The VexiiRiscv Implementation}

VexiiRiscv is a configurable RISC-V core implemented in SpinalHDL,
a Scala-based hardware description language.
Unlike traditional Verilog or VHDL, SpinalHDL supports
generative and highly parameterized hardware construction.
This approach allows for scalable design-space exploration, easier feature composition,
and cleaner abstraction boundaries. \cite{SpinalHDL2018, VexiiRiscVGithub}

The processor follows a classical pipelined architecture.
The VexiiRiscv processor supports up to 5 stages,
\textbf{Fetch}, \textbf{Decode}, \textbf{Execute}, \textbf{Memory} and \textbf{Write-back},
with \textbf{Fetch + Execute} (possibly unified) and \textbf{Memory and Write-back} (also unified).

One of the distinguishing characteristics of VexiiRiscv is its plugin-based architecture.
Major features—including memory management units (MMU), caches, branch predictors,
and custom execution units—are implemented as independent plugins.
The CXU extension is an execute extension. All computation done by CXU coprocessors is done
in the execution step of \textbf{VexiiRiscv}.

\section{The litex SoC Ecosystem}

\subsection{litex Framework Overview}

litex \cite{LiteXGitHub} is a python-based SoC construction framework built on top of Migen \cite{MiSoC2014}.
It enables automated generation of FPGA-based SoCs by integrating CPU cores,
memory controllers, interconnects, and peripherals into a cohesive hardware platform.

Within this ecosystem, litex acts as the integration layer between the VexiiRiscv CPU
core and FPGA peripherals such as UART, Ethernet \cite{LiteEthGithub}, timers \cite{LiteXGitHub},
and external memory controllers \cite{LiteDramGithub, LitePcieGithub}.
The framework abstracts low-level hardware connectivity and provides a unified build
system for generating bitstreams and associated firmware.

The framework automatically generates HDL, memory maps,  CSR definitions and software headers.

\subsection{Memory-Mapped Architecture}

Litex uses memory-mapped I/O.
Peripherals are accessed by reading or writing specific addresses:

\begin{verbatim}
#define UART_BASE 0xF0001000
\end{verbatim}

Software interacts with hardware by manipulating these addresses.

\begin{figure}[htb]
    \centering
    \begin{tikzpicture}[
        node distance=1.5cm and 0.5cm,
        master/.style={draw, rectangle, minimum width=3.5cm, minimum height=1.5cm, align=center, fill=red!10, rounded corners, font=\bfseries},
        slave/.style={draw, rectangle, minimum width=2.5cm, minimum height=1.5cm, align=center, fill=blue!10, rounded corners},
        bus/.style={draw, rectangle, minimum width=15cm, minimum height=1cm, align=center, fill=gray!20, font=\bfseries\large},
        arrow/.style={-Latex, thick},
        bidir/.style={Latex-Latex, thick},
        group/.style={draw, dashed, thick, fill=gray!5, rounded corners, inner sep=0.5cm}
        box/.style={rectangle, draw, rounded corners=3pt,
            minimum width=3.2cm, minimum height=0.75cm,
            text centered, font=\small}, 
    ]

    \node[bus] (wishbone) {Wishbone Interconnect (32-bit)};

    \node[master, above=1.5cm of wishbone, xshift=-4cm] (cpu) {VexiiRiscv CPU \\ \small + CXU Plugin};
    \node[master, above=1.5cm of wishbone, xshift=4cm, fill=orange!10] (dma) {LiteEth DMA \\ \small (Network Master)}; 

    \node[slave, below=1.5cm of wishbone, xshift=-6cm] (rom) {BootROM / \\ \small SPI Flash};
    \node[slave, below=1.5cm of wishbone, xshift=-3cm] (litedram) {LiteDRAM \\ \small (Main Memory)};
    \node[slave, below=1.5cm of wishbone, xshift=0cm] (liteuart) {LiteUART \\ \small (Serial Console)};
    \node[slave, below=1.5cm of wishbone, xshift=3cm] (liteeth) {LiteEth MAC \\ \small (Network Slave)};
    \node[slave, below=1.5cm of wishbone, xshift=6cm] (timer) {Timer / \\ \small Core CSRs};

    \draw[bidir] (cpu.south) -- (cpu.south |- wishbone.north) node[midway, left, font=\small] {Instr / Data};
    \draw[bidir] (dma.south) -- (dma.south |- wishbone.north) node[midway, right, font=\small] {RX / TX};

    \draw[bidir] (rom.north |- wishbone.south) -- (rom.north);
    \draw[bidir] (litedram.north |- wishbone.south) -- (litedram.north);
    \draw[bidir] (liteuart.north |- wishbone.south) -- (liteuart.north);
    \draw[bidir] (liteeth.north |- wishbone.south) -- (liteeth.north);
    \draw[bidir] (timer.north |- wishbone.south) -- (timer.north);

    \end{tikzpicture}
    \caption{Hardware architecture of the generated LiteX SoC}
\end{figure}

\subsection{Interconnect and Control}

litex commonly uses the Wishbone bus as its on-chip interconnect standard.
Wishbone is a simple and flexible synchronous bus protocol that supports
master-slave communication between cores and peripherals.
It enables memory-mapped I/O and simplifies peripheral integration. \cite{LiteXGitHub}

Control and Status Registers (CSRs) provide a structured mechanism for software-hardware communication.
CSRs are memory-mapped registers that allow software to configure hardware blocks,
query status information, and trigger operations \cite{LiteDramGithub, LiteXGitHub}.
This abstraction is essential for interacting with accelerators and managing runtime behavior.

\subsection{Memory Management}

litex integrates external memory controllers through LiteDRAM \cite{LiteDramGithub},
which provides support for SDRAM and DDR memory devices.
LiteDRAM abstracts memory timing, initialization, and refresh management,
exposing a consistent interface to the system bus.

The SoC memory map defines address regions for ROM, SRAM, DRAM, CSRs, and peripherals.
Proper memory layout is particularly important for Linux-based systems,
where kernel placement, device tree configuration,
and bootloader operation depend on predictable address organization.

\section{Software Development for litex-Based SoCs}

Developing software for a \textbf{litex}-based SoC requires
cross-compilation and careful linking against the target memory map.
Since litex systems typically use RISC-V softcores (e.g., VexRiscv),
programs must be compiled for the appropriate architecture. \cite{LiteXGitHub}

\subsection{Cross-Compilation}

Software is compiled on a host machine using a toolchain for the desired risc-v architecrure.

The resulting ELF binary is then loaded into on-chip memory
(SRAM, DRAM, or SPI Flash), as in \cite{CXUPlayground}.

\subsection{Linker Scripts}

litex SoCs define a memory map that must be respected during linking.
A linker script specifies instruction mrmory region, data memory region and stack and heap placement.

Correct linking ensures the program executes from valid memory
and interacts properly with memory-mapped peripherals.

\subsection{Bare-Metal vs Linux-Based Systems}

The Litex platform supports bare-metal, RTOS-based and linux operating systems.
Bare-metal software directly accesses memory-mapped registers,
while Linux- and RTOS-based systems require drivers and standard system calls.

\section{Buildroot}

\textbf{Buildroot} is an embedded Linux build system used to
generate complete root filesystems, toolchains, and kernel images
for embedded targets.

Buildroot automatescross-toolchain generation, linux kernel config,
rootfs and pakcage selection.
The workflow is as such: configuring the target,
selecting the architecture,
enabling desired packages and building the system complete image, while
the output typically includes the kernel image, DTB and the filesystem image.

Buildroot is commonly used with litex SoCs to produce
bootable Linux images for FPGA-based systems. It's used with \cite{CXUPlayground}.

\section{Custom Extension Units (CXU) Mechanism}

\subsection{Coupling Architecture}

Hardware accelerators can be integrated into a system using either
tightly coupled or loosely coupled approaches. \cite{Gray2022CFU, Gray2025CXU}

\textbf{Tightly coupled accelerators} are directly integrated into the CPU pipeline.
They operate as execution units and are invoked through custom instructions.
This approach offers low latency and direct register-file access but requires careful
hazard and pipeline management.

\textbf{Loosely coupled accelerators} are attached via system buses or DMA engines.
They operate as independent hardware blocks, often accessed through memory-mapped interfaces.
While simpler to integrate, they typically incur higher invocation latency.

The CXU mechanism adopts a tightly coupled strategy,
enabling instruction-level acceleration with minimal software overhead.

\begin{figure}[htb]
    \centering
    \begin{tikzpicture}[
        block/.style={draw, rectangle, minimum width=3cm, minimum height=4.5cm, align=center},
        arrow/.style={-Latex, thick},
        box/.style={rectangle, draw, rounded corners=3pt,
            minimum width=3.2cm, minimum height=0.75cm,
            text centered, font=\small}
    ]
        \node[block, fill=blue!10] (cpu) {VexiiRiscv CPU};
        \node[block, fill=green!10, right=4cm of cpu] (cxu) {CXU Accelerator};
        
        \draw[arrow] ([yshift=1.5cm]cpu.east) -- node[above, font=\small] {rs1 (32-bit)} ([yshift=1.5cm]cxu.west);
        \draw[arrow] ([yshift=0.5cm]cpu.east) -- node[above, font=\small] {rs2 (32-bit)} ([yshift=0.5cm]cxu.west);
        \draw[arrow] ([yshift=-0.5cm]cpu.east) -- node[above, font=\small] {cxidx / opcode} ([yshift=-0.5cm]cxu.west);
        
        \draw[arrow] ([yshift=-1.5cm]cxu.west) -- node[above, font=\small] {rd (32-bit)} ([yshift=-1.5cm]cpu.east);
        
        \draw[arrow, dashed] ([yshift=2.2cm]cpu.east) -- node[above, font=\small\itshape] {Valid / Cmd Enable} ([yshift=2.2cm]cxu.west);
        \draw[arrow, dashed] ([yshift=-2.2cm]cxu.west) -- node[above, font=\small\itshape] {Ready / Pipeline Stall} ([yshift=-2.2cm]cpu.east);
    \end{tikzpicture}
    \caption{CXU-CPU connection}
\end{figure}

\subsection{VexiiRiscv CXU User Plugin Interface}

The VexiiRiscv plugin interface enables custom instruction
integration at multiple levels of the pipeline.

\subsubsection{Instruction Decoding}

Specific opcode patterns are
reserved and mapped to the CXU \cite{SpinalHDL2018}.
During the decode stage, matching instructions are redirected to the custom execution logic by
SpinalHDL's built-in extension logic, after which the CXU connector extension sends the decoded
command to the CXU coprocessor.

\begin{figure}[htb]
    \centering
    \begin{enumerate}
        \item CXU state is updated through the use of $cxidx$ and $cxdata$ CSRs.
        \item The CXU custom command is called with $rs1$ and $rs2$ as inputs and $rd$ as output.
        \item Source operands are read from registers $rs1$ and $rs2$.
        \item The CXU processes the input state and registers and returns the desired output.
        \item The result is written back to the destination register $rd$.
        \item Updated CXU state may be read from through the use of $cxidx$ and $cxdata$ CSRs.
    \end{enumerate}
    \caption{VexiiRiscv data flow}
\end{figure}

Multi-cycle CXU operations require pipeline stall mechanisms
to preserve correctness. The plugin must signal readiness and manage structural,
data, and control hazards to ensure correct execution ordering.

\section{Linux Support on litex}

\subsection{Hardware Requirements}

Running Linux on a RISC-V SoC requires \cite{LinuxOnLiteXGitHub}
A Memory Management Unit (MMU) for virtual memory,
atomic instruction support (RV32A/RV64A) for synchronization primitives and
a supervisor privilege mode.

These features are enabled through appropriate VexiiRiscv plugins
and litex SoC configuration parameters.

\subsection{The Boot Process}

The Linux boot flow typically involves multiple stages.
OpenSBI \cite{OpenSBIGithub} provides the Supervisor Binary Interface (SBI),
acting as a firmware layer between machine mode and the operating system.
It handles low-level initialization and exposes standardized services to the kernel.

A secondary bootloader such as U-Boot \cite{UBootWebsite} loads the Linux kernel image and optional
initial ramdisk into memory. Once control is transferred,
the Linux kernel initializes drivers, memory management, and user-space services.

\begin{figure}[htb]
    \centering
    \begin{tikzpicture}[
        node distance=0.8cm,
        box/.style={draw, rectangle, minimum width=5cm, minimum height=1cm, rounded corners},
        box/.style={rectangle, draw, rounded corners=3pt,
            minimum width=3.2cm, minimum height=0.75cm, align=center,
            text centered, font=\small}
    ]
        \node[box, fill=gray!20] (reset) {hardware reset};
        \node[box, fill=blue!20, below=of reset] (litex) {litex bios (bootrom)\\ \small hardware init \& memory training};
        \node[box, fill=purple!20, below=of litex] (opensbi) {opensbi\\ \small machine mode environment setup};
        \node[box, fill=orange!20, below=of opensbi] (linux) {linux kernel\\ \small supervisor mode (mmu, dtb parsing)};
        \node[box, fill=cyan!20, below=of linux] (user) {buildroot user space\\ \small init scripts \& cxu applications};
        
        \draw[->] (reset) -- (litex);
        \draw[->] (litex) -- (opensbi);
        \draw[->] (opensbi) -- (linux);
        \draw[->] (linux) -- (user);
    \end{tikzpicture}
    \caption{the litex soc linux boot process}
\end{figure}

\subsection{Device Tree Overlay (DTS)}

The device tree describes the hardware layout to the Linux kernel \cite{DeviceTreeDoc}.
litex generates a hardware description that is translated into a Device Tree Source (DTS) file.
This file specifies CPU configuration, memory regions, peripheral addresses,
and interrupt mappings. The kernel uses this information to configure drivers
and manage hardware resources dynamically.

\section{Lattice-Based Cryptography}

\subsection{Post-Quantum Context}

Classical public-key cryptosystems such as RSA and elliptic curve cryptography
rely on number-theoretic problems vulnerable to Shor's algorithm on sufficiently powerful
quantum computers. This motivates the development of post-quantum cryptographic schemes.
\cite{Vranken2023UniformMul, Levi2020RISQV, Sozbilir2020Kyber,, Franchetti2024MultiWord}

Lattice-based cryptography is among the leading post-quantum approaches and forms the
basis of NIST-selected standards such as ML-KEM (formerly Kyber) and ML-DSA (formerly Dilithium).
These schemes rely on hardness assumptions related to structured lattice problems.

\section{The AES Encryption}

\textbf{AES} (Advanced Encryption Standard) \cite{NISTAES} is a symmetric encryption algorithm used to secure sensitive information. It was established by the U.S. National Institute of Standards and Technology and is widely used for encrypting sensitive information.
Due to it's definition, it m aps well onto hardware.

AES \cite{NISTAES} is based on a substitution-permutation network and can be efficiently implemented in both software and hardware. AES operates on a block of 16 bytes arranged in column-major order, called the \textbf{State}.

The key size used for an AES cipher specifies the number of transformation rounds that convert the plaintext into the ciphertext. The number of rounds are:

\begin{itemize}
    \item 10 rounds with 128-bit keys
    \item 12 rounds with 192-bit keys
    \item 14 rounds with 256-bit keys
\end{itemize}

Each round consists of several processing steps. A set of inverse rounds are applied during decryption to recover the plaintext.

The transformation steps are as such:

\begin{figure}[htb]
    \centering
    \begin{tikzpicture}[
        node distance=0.6cm and 1.5cm,
        box/.style={draw, rectangle, minimum width=3.5cm, minimum height=0.8cm, align=center, fill=blue!10, rounded corners},
        keybox/.style={draw, rectangle, minimum width=2.5cm, minimum height=0.8cm, align=center, fill=red!10, rounded corners},
        io/.style={draw, trapezium, trapezium left angle=70, trapezium right angle=110, minimum width=3cm, minimum height=0.8cm, align=center, fill=green!10},
        arrow/.style={-Latex, thick},
        group/.style={draw, rectangle, dashed, thick, inner sep=0.4cm, fill=gray!10, rounded corners}
    ]
    \node[io] (plaintext) {Plaintext (128-bit)};
    
    \node[box, below=1cm of plaintext] (init_addkey) {AddRoundKey};
    \node[keybox, right=of init_addkey] (key0) {Round Key 0};
    
    \node[box, below=1.2cm of init_addkey] (subbytes) {SubBytes};
    \node[box, below=of subbytes] (shiftrows) {ShiftRows};
    \node[box, below=of shiftrows] (mixcols) {MixColumns};
    \node[box, below=of mixcols] (main_addkey) {AddRoundKey};
    \node[keybox, right=of main_addkey] (key1_9) {Round Keys 1--9};

    \node[box, below=2cm of main_addkey] (subbytes_f) {SubBytes};
    \node[box, below=of subbytes_f] (shiftrows_f) {ShiftRows};
    \node[box, below=of shiftrows_f] (fin_addkey) {AddRoundKey};
    \node[keybox, right=of fin_addkey] (key10) {Round Key 10};

    \node[io, below=1cm of fin_addkey] (ciphertext) {Ciphertext (128-bit)};

    \begin{scope}[on background layer]
        \node[group, fit=(subbytes) (shiftrows) (mixcols) (main_addkey), 
              label=left:\rotatebox{90}{\textbf{9 Standard Rounds}}] (main_loop) {};
              
        \node[group, fit=(subbytes_f) (shiftrows_f) (fin_addkey), 
              label=left:\rotatebox{90}{\textbf{Final Round}}] (final_loop) {};
    \end{scope}

    \draw[arrow] (plaintext) -- (init_addkey);
    \draw[arrow] (key0) -- (init_addkey);
    
    \draw[arrow] (init_addkey) -- (subbytes);
    
    \draw[arrow] (subbytes) -- (shiftrows);
    \draw[arrow] (shiftrows) -- (mixcols);
    \draw[arrow] (mixcols) -- (main_addkey);
    \draw[arrow] (key1_9) -- (main_addkey);

    \draw[arrow, rounded corners=5pt] (main_addkey.south) -- ++(0,-0.4) -| 
          ([xshift=-1cm]main_loop.west) |- node[pos=0.25, left] {Loop 9x} 
          ([yshift=0.4cm]subbytes.north) -- (subbytes.north);
    
    \draw[arrow] (main_addkey) -- (subbytes_f);
    
    \draw[arrow] (subbytes_f) -- (shiftrows_f);
    \draw[arrow] (shiftrows_f) -- (fin_addkey);
    \draw[arrow] (key10) -- (fin_addkey);

    \draw[arrow] (fin_addkey) -- (ciphertext);

    \end{tikzpicture}
    \caption{Flowchart of the AES encryption algorithm}
\end{figure}

\subsection{AES Round Operations}

\subsubsection{KeyExpansion}

The KeyExpansion step derives round keys from the original cipher key.

\subsubsection{AddRoundKey}

AddRoundKey performs a byte-wise XOR between the State and the round key.

\subsubsection{MixColumns}

MixColumns mixes each column using matrix multiplication in $GF(2^8)$.

Each column is multiplied by the fixed matrix:

$$
\begin{bmatrix}
02 & 03 & 01 & 01 \\
01 & 02 & 03 & 01 \\
01 & 01 & 02 & 03 \\
03 & 01 & 01 & 02
\end{bmatrix}
$$

\subsection{Speeding up AES with the CXU}

AES could be sped up by having the arithmetic operations in $GF(2^8)$ be done by the extension.
Instead of doing a loop of 5-6 iterations and doing 3-4 arithmetic operations in each iteration,
an extension may be created to execute the arithmetic oepration in a finite field in a single CPU cycle.

\section{GZIP}

\textbf{GZIP} (GNU Zip) \cite{GZIPWebsite} is a widely used lossless compression format.
It is based on the \textbf{DEFLATE} compression algorithm \cite{DEFLATESpec}, which combines
dictionary-based compression and entropy encoding to reduce data size efficiently.

A GZIP file consists of:

\begin{itemize}
    \item A header (metadata, flags, optional filename)
    \item Compressed data (DEFLATE stream)
    \item A footer (CRC-32 and original size)
\end{itemize}

\subsection{The DEFLATE Algorithm}

The DEFLATE algorithm combines the LZ77 compression and huffman coding.
The compression pipeline is therefore
\begin{enumerate}
    \item To detect repeated sequences (LZ77)
    \item To encode literals and length-distance pairs
    \item To apply Huffman coding to the symbol stream
    \item To pack the result into a bit-stream
\end{enumerate}

\begin{figure}[htb]
\centering
\begin{tikzpicture}[
    node distance=0.7cm and 1.5cm,
    box/.style={draw, rectangle, minimum width=3.5cm, minimum height=0.8cm, align=center, fill=blue!10, rounded corners},
    keybox/.style={draw, rectangle, minimum width=2.5cm, minimum height=0.8cm, align=center, fill=red!10, rounded corners},
    io/.style={draw, trapezium, trapezium left angle=70, trapezium right angle=110, minimum width=3cm, minimum height=0.8cm, align=center, fill=green!10},
    arrow/.style={-Latex, thick},
    group/.style={draw, rectangle, dashed, thick, inner sep=0.4cm, fill=gray!10, rounded corners}
]

\node[io] (input) {Input Data Stream};

\node[box, below=1cm of input] (lz77) {LZ77 Sliding Window\\Match Detection};

\node[box, below=of lz77] (tokens) {Generate Symbols\\(Literals, Length, Distance)};

\node[box, below=1.2cm of tokens] (freq) {Symbol Frequency\\Counting};

\node[keybox, right=of freq] (tree) {Huffman Tree\\Generation};

\node[box, below=of freq] (huff) {Huffman Encoding};

\node[box, below=of huff] (pack) {Bitstream Packing};

\node[io, below=1cm of pack] (output) {Compressed DEFLATE Stream};

\begin{scope}[on background layer]

\node[group, fit=(lz77) (tokens),
      label=left:\rotatebox{90}{\textbf{LZ77 Stage}}] (lzgroup) {};

\node[group, fit=(freq) (huff),
      label=left:\rotatebox{90}{\textbf{Huffman Coding}}] (huffgroup) {};

\end{scope}

\draw[arrow] (input) -- (lz77);
\draw[arrow] (lz77) -- (tokens);

\draw[arrow] (tokens) -- (freq);

\draw[arrow] (freq) -- (huff);
\draw[arrow] (tree) -- (huff);

\draw[arrow] (huff) -- (pack);
\draw[arrow] (pack) -- (output);

\end{tikzpicture}
\caption{Flowchart of the DEFLATE compression pipeline}
\end{figure}

\subsection{LZ77 Compression}

LZ77 operates using a sliding window over previously seen data.
Repeated sequences are replaced with \textit{(length, distance)} pairs.

After LZ77 processing, the data stream consists of:

\begin{enumerate}
    \item Literal bytes (0--255)
    \item Length codes (representing match length)
    \item Distance codes (representing match distance)
\end{enumerate}

The symbols are not stored directly; they are entropy encoded using Huffman coding.

\subsection{Huffman Coding}

Huffman coding assigns shorter bit patterns to frequently occurring symbols.

\subsection{Dynamic vs Static Huffman Trees}

DEFLATE supports:

\begin{itemize}
    \item \textbf{Static Huffman}: predefined code tables
    \item \textbf{Dynamic Huffman}: code tables computed per block
\end{itemize}

Dynamic coding usually achieves better compression because it adapts
to the local symbol frequency of each block.

Such characteristics make DEFLATE a strong candidate for
instruction-level acceleration via the CXU mechanism,
particularly for:

\begin{itemize}
    \item Fast match search primitives
    \item Parallel prefix bit-packing
    \item Hardware-assisted CRC computation
\end{itemize}

In contrast to AES, which relies on algebraic transformations in $GF(2^8)$,
GZIP compression primarily relies on pattern detection and entropy encoding
to eliminate statistical redundancy in data.

\subsection{Custom Instructions for DEFLATE Acceleration}

To accelerate critical primitives in the DEFLATE pipeline,
several custom instructions were implemented using
the CXU extension mechanism.
These instructions target operations that occur frequently in
GZIP compression and decompression, particularly bit manipulation,
Huffman decoding support, and address calculations for LZ77 back-references.

The implemented instructions are described below.

\subsubsection{Bit Reversal}

The \textbf{cxu\_bitrev} instruction reverses the order of bits in a
32-bit word. Bit reversal is useful in several bitstream manipulation
tasks, including efficient construction of lookup indices for
variable-length codes.

Given an input value $x$, the instruction produces an output where
bit $i$ of the result corresponds to bit $31-i$ of the input.

$$
y_i = x_{31-i}
$$

This operation is implemented purely combinationally by wiring the
input bits to reversed output positions.

\subsubsection{Bit Mask Generation}

The \textbf{cxu\_mask} instruction generates a mask containing the
lowest $n$ bits set to one.

$$
mask = (1 \ll n) - 1
$$

Such masks are frequently required when extracting variable-length
bit fields from the compressed DEFLATE bitstream. Generating them
in a single instruction reduces the number of shifts and arithmetic
operations required in software.

\subsubsection{Huffman Table Index Extraction}

The \textbf{cxu\_huft\_idx} instruction computes a masked index value:

$$
index = x \;\&\; ((1 \ll n) - 1)
$$

This operation is commonly used during Huffman decoding when
determining the lookup index into a decoding table based on the
current bit buffer. By fusing the mask generation and bit extraction
into a single instruction, the number of required operations in the
decoding loop is reduced.

\subsubsection{LZ77 Back-Reference Address Calculation}

The \textbf{cxu\_copy\_addr} instruction accelerates address
calculation for LZ77 back-references during decompression.

In the DEFLATE algorithm, repeated sequences are encoded as
\textit{(length, distance)} pairs. During decompression, the
distance value determines how far back in the output window the
copy operation should read.

The instruction performs the following computation:

$$
addr = (w - dist\_base - extra) \bmod 2^{15}
$$

where:

\begin{itemize}
\item $w$ is the current write position in the sliding window,
\item $dist\_base$ is the base distance value from the Huffman code,
\item $extra$ represents additional bits extending the distance.
\end{itemize}

By combining subtraction and masking into a single hardware
operation, the instruction reduces the arithmetic overhead in the
inner decompression loop.

\subsubsection{Implementation Characteristics}

The CXU module is entirely combinational and does not maintain
internal state. All instructions operate on two 32-bit operands
and produce a single 32-bit result in the same cycle.

Because no state is written, the state interface of the CXU unit
remains unused. This simplifies integration and ensures that the
custom instructions behave similarly to standard arithmetic
instructions from the perspective of the processor pipeline.

These instructions target frequently executed primitives in the
DEFLATE decoding path and therefore provide opportunities for
improving overall decompression throughput.

